\section{Kesimpulan dan Rekomendasi}
\label{sec:kesimpulan}

\subsection{Kesimpulan}

PLTS atap dan pensiun dini PLTU belum berjalan bersamaan karena keduanya diatur terpisah, pemilihan PLTU belum menimbang jenis kontrak, kompensasi dan pengganti tidak dipersyaratkan sebelum penutupan, dan tidak ada alat hitung bersama. Program pasangan kuota PLTS atap dan pensiun dini menjawab celah tersebut dengan memasangkan setiap PLTU yang dipensiunkan dengan tambahan PLTS atap setara, mendahulukan PLTU IPP, dan menetapkan syarat penutupan yang jelas.

Pada sistem Jawa-Bali, roadmap Fase 1-3 memensiunkan 3,0 GW PLTU IPP, menambah 3,19 GWp PLTS atap, dan menghindari 9.572 kematian dini serta 35,4 juta ton CO$_2$ selama 2026-2040. Beban negara bersih Rp11,9 triliun, kurang dari Rp1 triliun per tahun, dengan dana luar Rp15,2 triliun yang masih di bawah alokasi CIPP JETP. Pensiun dini PLTU IPP adalah bagian yang membuat program terjangkau, karena PLTS atap tanpa pensiun dini justru membebani negara lebih besar. Kelayakan ini bergantung pada besaran kompensasi, sehingga kesepakatan kompensasi dan pengikatan dana luar harus mendahului keputusan penutupan.

\subsection{Rekomendasi}

\textbf{Kepada Menteri Energi dan Sumber Daya Mineral:}
\begin{enumerate}
  \item Menetapkan peraturan kuota PLTS atap berpasangan dengan rencana pensiun dini (rancangan pada Lampiran~\ref{lamp:permen}) pada Fase 1.
  \item Menambahkan jenis kontrak dan beban kesehatan per lokasi sebagai pertimbangan dalam penilaian kriteria Permen ESDM 10/2025.
  \item Menetapkan kewajiban PLTS atap gedung instansi pemerintah di Jawa-Bali mulai 2028.
\end{enumerate}

\textbf{Kepada Menteri Keuangan:}
\begin{enumerate}[resume]
  \item Menyediakan insentif fiskal 15\% bagi PLTS atap industri dan anggaran PLTS atap gedung pemerintah.
  \item Mensyaratkan pengikatan paling sedikit 80\% kompensasi dari dana luar sebelum keputusan penutupan.
\end{enumerate}

\textbf{Kepada Direktur Utama PT PLN (Persero):}
\begin{enumerate}[resume]
  \item Memasukkan data kontrak IPP ke PowerShift untuk penyaringan kandidat dan membuka hasil agregatnya.
  \item Menjadwalkan pembangkit pengganti dalam RUPTL terhadap tanggal penutupan Cirebon-1, Tanjung Jati B 1-2, Cilacap Phase I, dan Celukan Bawang.
  \item Memperkuat jaringan distribusi Jawa-Bali untuk menyerap 3 GWp PLTS atap tambahan.
\end{enumerate}

\textbf{Kepada Satuan Tugas Transisi Energi:}
\begin{enumerate}[resume]
  \item Memakai PowerShift sebagai alat hitung bersama untuk penetapan daftar urut, pemeriksaan syarat penutupan, dan tinjauan indikator di akhir setiap fase.
  \item Melaporkan indikator kesetaraan daya bersama indikator energi agar substitusi tidak dinilai berlebihan.
\end{enumerate}

\textbf{Kepada Kementerian Ketenagakerjaan dan pemerintah daerah:}
\begin{enumerate}[resume]
  \item Memulai pelatihan ulang pekerja setiap PLTU kandidat paling lambat satu tahun sebelum penutupan, diarahkan ke pemasangan dan pemeliharaan PLTS atap.
\end{enumerate}

Dengan langkah-langkah tersebut, PLTS atap tidak lagi dipandang sebagai pengurang pendapatan PLN, dan pensiun dini tidak lagi dipandang sebagai beban kompensasi semata. Keduanya menjadi satu program yang biayanya terukur dan manfaatnya dapat dibuktikan.
